\heading{实验物理中的统计方法-统计误差、置信区间和极限}





\begin{question}
假设估计量 \(\hat\theta\) 服从高斯分布，高斯分布的参数分别为
\(\hat\theta\) 的真值 \(\theta\) 和标准偏差
\(\sigma_{\hat\theta}\)。假设 \(\sigma_{\hat\theta}\) 已知。

\begin{enumerate}
\item
  画出定义置信带的函数 \(u_\alpha(\theta)\) 和 \(v_\beta(\theta)\)（参见
  Statistical Data Analysis 第 9.2 节）。
\item
  证明置信水平为 \(1-\gamma\) 时参数 \(\theta\) 的中心置信区间由下式给出
\end{enumerate}

\[
\left[\hat\theta-\sigma_{\hat\theta}\Phi^{-1}(1-\gamma/2),\ \hat\theta+\sigma_{\hat\theta}\Phi^{-1}(1-\gamma/2)\right],
\]

其中 \(\Phi^{-1}\) 是标准高斯分布的分位数。
\begin{solution}
设估计量 \(\hat\theta\) 服从高斯分布，均值为真值
\(\theta\)，标准差为已知常数 \(\sigma_{\hat\theta}\)。

\begin{enumerate}
\item 置信带函数

若以 \(1-\alpha\) 和 \(1-\beta\) 作为上下边界的置信带，则定义

\[
P(\hat\theta\le u_\alpha(\theta)\mid \theta)=1-\alpha,
\qquad
P(\hat\theta\le v_\beta(\theta)\mid \theta)=\beta.
\]

由于

\[
\frac{\hat\theta-\theta}{\sigma_{\hat\theta}}\sim N(0,1),
\]

故

\[
u_\alpha(\theta)=\theta+\sigma_{\hat\theta}\Phi^{-1}(1-\alpha),
\qquad
v_\beta(\theta)=\theta+\sigma_{\hat\theta}\Phi^{-1}(\beta).
\]

\item 中心置信区间

对中心区间取 \(\alpha=\beta=\gamma/2\)，反演置信带后得到

\[
\theta\in
\Big[
\hat\theta-\sigma_{\hat\theta}\Phi^{-1}(1-\gamma/2),
\hat\theta+\sigma_{\hat\theta}\Phi^{-1}(1-\gamma/2)
\Big].
\]

这就是式 (9.1)。
\pyfig[0.72\linewidth]{figures/fig_09_01_gaussian_belt.png}{已知标准差的高斯估计量置信带；反演横向截线得到参数区间。}\end{enumerate}
\end{solution}
\end{question}



\begin{question}
随机变量 \(x\) 服从均值为 \(\xi\) 的指数分布，考虑对 \(x\) 的 \(n\)
次观测。参数 \(\xi\) 的最大似然估计量（见 Statistical Data Analysis
(6.6) 式）由下式给出

\[
\hat\xi=\frac1n\sum_{i=1}^{n}x_i,
\]

并且 \(\hat\xi\) 的概率密度函数（参见 Statistical Data Analysis (10.25)
式）为

\[
g(\hat\xi;\xi)=\frac{n^n}{(n-1)!}\frac{\hat\xi^{n-1}}{\xi^n}e^{-n\hat\xi/\xi}.
\]

\begin{enumerate}
\item
  证明：定义置信带的曲线 \(u_\alpha(\xi)\) 和 \(v_\beta(\xi)\) 为
\end{enumerate}

\[
\begin{split}
u_\alpha(\xi)&=\frac{\xi}{2n}F_{\chi^2}^{-1}(1-\alpha;2n),\\
v_\beta(\xi)&=\frac{\xi}{2n}F_{\chi^2}^{-1}(\beta;2n),
\end{split}
\]

其中 \(F_{\chi^2}^{-1}\) 为 \(\chi^2\) 分布的分位数。根据习题
2.6，\(\chi^2\) 分布的累积分布可以与不完全伽马函数 \(P(x,n)\) 联系起来：

\[
F_{\chi^2}(2x;2n)=P(x,n)\equiv\frac1{\Gamma(n)}\int_0^x e^{-t}t^{n-1}\,dt.
\]

取 \(\alpha=\beta=0.159\)，\(n=5\)，画出 \(u_\alpha(\xi)\) 和
\(v_\beta(\xi)\)。（\(\chi^2\)
分布的分位数可以从标准分布表中查出，或者在 ROOT 中调用
\texttt{TMath::ChisquareQuantile(Double\_t\ p,\ Double\_t\ ndf)}
函数得到。）

\begin{enumerate}[start=2]
\item
  求出置信区间 \([a,b]\) 作为估计值 \(\hat\xi\)、样本容量 \(n\)
  以及置信水平 \(\alpha\) 和 \(\beta\) 的函数。假设估计值为
  \(\hat\xi=1.0\)，在 \(u_\alpha(\xi)\) 和 \(v_\beta(\xi)\)
  的图上画出该估计量的值。取 \(n=5\)，\(\alpha=\beta=0.159\)，计算 \(a\)
  和 \(b\)。将计算结果与估计值加减一倍标准差得到的区间进行比较。
\end{enumerate}
\begin{solution}
设 \(x\) 服从均值为 \(\xi\) 的指数分布，样本均值的 MLE 为

\[
\hat\xi=\frac1n\sum_{i=1}^n x_i.
\]

其密度为

\[
g(\hat\xi;\xi)=\frac{n^n}{(n-1)!}\frac{\hat\xi^{n-1}}{\xi^n}e^{-n\hat\xi/\xi}.
\]

\begin{enumerate}
\item 置信带函数

注意到

\[
\frac{2n\hat\xi}{\xi}\sim\chi^2_{2n},
\]

故若定义

\[
P(\hat\xi\le u_\alpha(\xi)\mid\xi)=1-\alpha,
\qquad
P(\hat\xi\le v_\beta(\xi)\mid\xi)=\beta,
\]

则

\[
\frac{2n\,u_\alpha(\xi)}{\xi}=F_{\chi^2}^{-1}(1-\alpha;2n),
\qquad
\frac{2n\,v_\beta(\xi)}{\xi}=F_{\chi^2}^{-1}(\beta;2n).
\]

因此

\[
u_\alpha(\xi)=\frac{\xi}{2n}F_{\chi^2}^{-1}(1-\alpha;2n),
\qquad
v_\beta(\xi)=\frac{\xi}{2n}F_{\chi^2}^{-1}(\beta;2n).
\]

\item 反演区间

给定观测 \(\hat\xi\)，反演后得到

\[
a=\frac{2n\hat\xi}{F_{\chi^2}^{-1}(1-\alpha;2n)},
\qquad
b=\frac{2n\hat\xi}{F_{\chi^2}^{-1}(\beta;2n)}.
\]

题目给出 \(\hat\xi=1.0\)，\(n=5\)，\(\alpha=\beta=0.159\)，程序算得

\[
[a,b]=[0.69845,
1.75902].
\]

若直接用``大样本一倍标准差''近似，因

\[
\sigma_{\hat\xi}=\frac{\xi}{\sqrt n},
\]

常见近似区间为

\[
1\pm\frac{1}{\sqrt5}=[0.5528,1.4472].
\]

可以看出精确区间明显不对称，而高斯近似过于对称。
\pyfig[0.72\linewidth]{figures/fig_09_02_exponential_belt.png}{\(n=5\)、\(\alpha=\beta=0.159\) 时指数均值参数的置信带与 \(\hat\xi=1\) 的反演。}\end{enumerate}
\end{solution}
\end{question}



\begin{question}
证明二项分布的参数 \(p\) 的上限和下限为

\[
p_{\rm lo}=\frac{nF_F^{-1}[\alpha;2n,2(N-n+1)]}{N-n+1+nF_F^{-1}[\alpha;2n,2(N-n+1)]},
\]

\[
p_{\rm up}=\frac{(n+1)F_F^{-1}[1-\beta;2(n+1),2(N-n)]}{(N-n)+(n+1)F_F^{-1}[1-\beta;2(n+1),2(N-n)]}.
\]

其中上下限的置信水平分别为 \(1-\alpha\) 和 \(1-\beta\)，\(n\) 为 \(N\)
次试验中成功的次数，\(F_F^{-1}\) 为 \(F\) 分布的分位数，由 \(F\)
分布定义：

\[
f(x;n_1,n_2)=\left(\frac{n_1}{n_2}\right)^{n_1/2}
\frac{\Gamma\left(\frac12(n_1+n_2)\right)}{\Gamma\left(\frac12n_1\right)\Gamma\left(\frac12n_2\right)}
x^{n_1/2-1}\left(1+\frac{n_1}{n_2}x\right)^{-(n_1+n_2)/2},
\]

其中 \(x>0\)，参数 \(n_1\) 和 \(n_2\)
为整数（自由度）。利用二项分布累积分布函数与自由度为 \(n_1=2(n+1)\) 和
\(n_2=2(N-n)\) 的累积分布函数 \(F_F(x)\) 的关系

\[
\sum_{k=0}^{n}\frac{N!}{k!(N-k)!}p^k(1-p)^{N-k}
=1-F_F\left[\frac{(N-n)p}{(n+1)(1-p)};2(n+1),2(N-n)\right].
\]

\(F\) 分布的分位数可以从标准分布表中查得，或者调用 ROOT 中的函数计算。
\begin{solution}
二项分布参数 \(p\) 的精确区间来自对二项分布尾概率的反演。对观测到
\(n\) 个成功、总试验数为 \(N\) 的情形，下限 \(p_{\rm lo}\) 由
\[
P(X\ge n\mid p_{\rm lo})=1-\sum_{k=0}^{n-1}\binom Nk p_{\rm lo}^k(1-p_{\rm lo})^{N-k}=1-\alpha
\]
等价地给出；上限 \(p_{\rm up}\) 由
\[
P(X\le n\mid p_{\rm up})=1-\beta
\]
给出。利用二项尾概率与 \(F\) 分布累积分布之间的恒等式，可以把这两个反演方程写成
\(F\) 分布分位数的形式。

具体地，二项分布的 Clopper--Pearson 区间也可写成 Beta 分布分位数
\[
p_{\rm lo}=B^{-1}(\alpha;n,N-n+1),
\qquad
p_{\rm up}=B^{-1}(1-\beta;n+1,N-n),
\]
其中端点情形 \(n=0\) 或 \(n=N\) 需按极限理解。
利用 Beta 分布与 \(F\) 分布分位数的关系
\[
B^{-1}(q;a,b)=
\frac{aF_F^{-1}(q;2a,2b)}{b+aF_F^{-1}(q;2a,2b)},
\]
分别取 \((a,b)=(n,N-n+1)\) 和 \((a,b)=(n+1,N-n)\)，即得
\[
p_{\rm lo}=
\frac{nF_F^{-1}[\alpha;2n,2(N-n+1)]}
{N-n+1+nF_F^{-1}[\alpha;2n,2(N-n+1)]},
\]
\[
p_{\rm up}=
\frac{(n+1)F_F^{-1}[1-\beta;2(n+1),2(N-n)]}
{(N-n)+(n+1)F_F^{-1}[1-\beta;2(n+1),2(N-n)]}.
\]
这就是题目要求证明的上下限公式。
\end{solution}
\end{question}



\begin{question}
在 CERN 的 Gargamelle
气室中进行了反中微子-核子散射实验。选择的事例是只产生强子的事例（通过中性流过程
\(\bar\nu_\mu N\to\bar\nu_\mu X\)），或者是产生强子和一个 \(\mu\)
子的事例（通过带电流过程 \(\bar\nu_\mu N\to\mu^+X\)）。从 212
个事例样本中，64 个事例归类为中性流（NC）过程，148
个归类为带电流（CC）过程。估计事例为中性流（NC）的概率，并求出 68.3\%
的中心置信区间。求 NC 与 CC 过程概率之比的估计量和置信区间。
\begin{solution}
样本中共有 212 个事件，其中 64 个是 NC，148 个是 CC，因此

\[
\hat p=\frac{64}{212}=0.301887.
\]

按 68.3\% 中心精确区间（Clopper--Pearson）计算，程序得到

\[
p\in[0.26906,0.33672].
\]

NC 与 CC 概率之比为

\[
r=\frac{p}{1-p}.
\]

其 MLE 为

\[
\hat r=\frac{64}{148}=0.43243.
\]

因为 \(r(p)=p/(1-p)\) 是单调函数，区间可直接变换：

\[
r\in\left[
\frac{p_{\rm lo}}{1-p_{\rm lo}},
\frac{p_{\rm up}}{1-p_{\rm up}}
\right]
=[0.36810,0.50766].
\]
\end{solution}
\end{question}



\begin{question}
在研究粒子碰撞的实验中，选择出来 10 个某种类型的事例，比如说某个量 \(x\)
具有较高值的事例。这 10 个高 \(x\) 值的事例中，发现有 2 个事例包含
\(\mu\) 子。

\begin{enumerate}
\item
  高 \(x\) 值事例中包含 \(\mu\) 子的事例数目服从二项分布，求参数 \(p\)
  的 68.3\% 中心置信区间。将结果表示为 \(p=\hat p^{+d}_{-c}\)，其中
  \(\hat p\) 为 \(p\) 的最大似然估计，\([\hat p-c,\hat p+d]\)
  为置信区间。
\item
  将 (a) 中的区间与 \(\hat p\pm\hat\sigma_{\hat p}\) 进行比较，其中
  \(\hat\sigma_{\hat p}\) 为 \(\hat p\) 的标准偏差的估计量。
\item
  经常犯的错误是将高 \(x\) 值的事例数 10 当作随机变量，并将其方差引入
  \(\hat p\) 的误差中（例如通过误差传递）。为什么这种方法是不正确的？
\end{enumerate}
\begin{solution}
\begin{enumerate}
\item 精确中心区间

这里样本容量固定为 10，所以
\[
n\sim\operatorname{Binom}(10,p),
\qquad
\hat p=\frac{2}{10}=0.2.
\]
取 \(68.3\%\) 中心区间，即两侧尾概率各约为 \((1-0.683)/2=0.1585\)。
Clopper--Pearson 精确区间为
\[
p\in[0.07191,0.40555].
\]
写成 \(p=\hat p^{+d}_{-c}\)，有
\[
c=0.2-0.07191=0.12809,
\qquad
d=0.40555-0.2=0.20555.
\]

\item 与 \(\hat p\pm\hat\sigma_{\hat p}\) 比较

大样本近似为
\[
\hat\sigma_{\hat p}=\sqrt{\frac{\hat p(1-\hat p)}{N}}.
\]
代入 \(\hat p=0.2\)、\(N=10\)：
\[
\hat\sigma_{\hat p}=\sqrt{\frac{0.2\times0.8}{10}}=0.12649.
\]
故
\[
\hat p\pm\hat\sigma_{\hat p}=[0.07351,0.32649].
\]
精确区间明显不对称；简单 \(\pm1\sigma\) 近似把上误差低估得较多。

\item 为什么不能把 10 也当随机量传播误差

题目条件是“先选出 10 个高 \(x\) 事例”，再在这 10 个事例中数含 \(\mu\) 子的个数。也就是说，正确模型是条件在固定样本容量 \(N=10\) 下的二项分布。若再给 \(N=10\) 加上泊松或其它波动，就相当于改成了另一个无条件实验模型，会把“高 \(x\) 总率的不确定性”错误地混入“给定高 \(x\) 后含 \(\mu\) 的条件概率”中。
\end{enumerate}
\end{solution}
\end{question}



\begin{question}
假设为了产生习题 (9.5) 中的事例，收集的总数据对应于积分亮度为
\(L=1\,\mathrm{pb}^{-1}\)（误差可以忽略）。产生的给定类型事例总数可以看作均值为
\(\nu=\sigma L\) 的泊松随机变量 \(n\)，其中 \(\sigma\)
为产生截面。（为什么是泊松分布？）

\begin{enumerate}
\item
  对于高 \(x\) 值事例以及高 \(x\) 值的 \(\mu\)
  子事例，假设观测到的事例数分别为 \(n_x=10\) 和
  \(n_{x\mu}=2\)，求事例数期待值 \(\nu_x\) 和 \(\nu_{x\mu}\) 的 68.3\%
  中心置信区间。对应的产生截面 \(\sigma_x\) 和 \(\sigma_{x\mu}\)
  的置信区间为多少？
\item
  将 (a) 中得到的置信区间与通过加减一倍标准偏差得到的区间进行比较。
\item
  假设另外一个实验的积分亮度为 \(L'=100\,\mathrm{pb}^{-1}\)，观测到
  \(n_x'\) 个高 \(x\) 值的事例。但是这个实验不能鉴别 \(\mu\)
  子。利用数据 \(n_x\)，\(n_{x\mu}\) 和 \(n_x'\) 构造参数 \(\sigma_x\)
  和 \(\sigma_{x\mu}/\sigma_x\) 的对数似然函数。证明 \(p\)
  的最大似然估计量与 \(n_x'\) 无关。这是否说明第二个实验的结果对 \(p\)
  的估计没有影响？
\item
  假设最初的实验没有测量高 \(x\) 值的事例数，只是测得了包含 \(\mu\)
  子的高 \(x\) 值事例数。利用结果 \(n_{x\mu}=2\) 和 \(n_x'\)，构造
  \(\sigma_x\) 和 \(p\)
  的对数似然函数，并求出最大似然估计量。利用误差传递估计 \(\hat p\)
  的标准偏差，并比较区间 \(\hat p\pm\sigma_{\hat p}\) 与习题 (9.5) 中
\begin{enumerate}
\item 和 (b) 得到的区间。选作：这种情况下，如何构造 \(p\) 的置信区间？
\end{enumerate}
\end{enumerate}
\begin{solution}
设 \(L=1\,\mathrm{pb}^{-1}\)，观测到
\[
n_x=10,
\qquad
n_{x\mu}=2.
\]
事例计数来自稀有独立碰撞，在固定积分亮度下可视为泊松变量，因此
\[
n_x\sim\operatorname{Pois}(\nu_x),
\qquad
n_{x\mu}\sim\operatorname{Pois}(\nu_{x\mu}),
\]
其中 \(\nu_x=\sigma_xL\)、\(\nu_{x\mu}=\sigma_{x\mu}L\)。

\begin{enumerate}
\item 泊松均值和截面的精确区间

泊松均值 \(\nu\) 的中心区间可由 \(\chi^2\) 分位数给出：
\[
\nu_{\rm lo}=\frac12F_{\chi^2}^{-1}(\alpha/2;2n),
\qquad
\nu_{\rm up}=\frac12F_{\chi^2}^{-1}(1-\alpha/2;2(n+1)),
\]
其中 \(\alpha=1-0.683\)。代入得
\[
\nu_x\in[6.8897,14.2695],
\qquad
\nu_{x\mu}\in[0.7077,4.6394].
\]
因 \(L=1\,\mathrm{pb}^{-1}\)，所以数值上
\[
\sigma_x\in[6.8897,14.2695]\,\mathrm{pb},
\qquad
\sigma_{x\mu}\in[0.7077,4.6394]\,\mathrm{pb}.
\]

\item 与 \(\pm1\sigma\) 比较

泊松大样本近似给出
\[
10\pm\sqrt{10}=[6.84,13.16],
\qquad
2\pm\sqrt2=[0.586,3.414].
\]
对 \(n=10\) 尚可，对 \(n=2\) 则明显过于对称且上端偏低。

\item 第二个实验只测高 \(x\) 总数

若另一个实验亮度为 \(L'=100\,\mathrm{pb}^{-1}\)，观测到 \(n'_x\) 个高
\(x\) 事例且不能识别 \(\mu\)，则可写联合似然
\[
L(\sigma_x,p)=
\operatorname{Pois}(n_x\mid\sigma_xL)
\operatorname{Binom}(n_{x\mu}\mid n_x,p)
\operatorname{Pois}(n'_x\mid\sigma_xL'),
\]
其中 \(p=\sigma_{x\mu}/\sigma_x\)。与 \(p\) 有关的项只有
\[
n_{x\mu}\log p+(n_x-n_{x\mu})\log(1-p),
\]
故
\[
\hat p=\frac{n_{x\mu}}{n_x},
\]
与 \(n'_x\) 无关。在这个条件二项模型下，第二个实验确实不改变
\(p\) 的剖面似然形状；它主要改善 \(\sigma_x\) 的估计。若存在共同系统误差或采用不同的无条件参数化，结论才可能改变。

\item 第一个实验只测到 \(n_{x\mu}\)

若原实验没有测 \(n_x\)，只有 \(n_{x\mu}\)，再结合第二个实验的 \(n'_x\)，则联合似然为
\[
L(\sigma_x,p)=
\operatorname{Pois}(n_{x\mu}\mid p\sigma_xL)
\operatorname{Pois}(n'_x\mid\sigma_xL').
\]
对数似然为
\[
\ell=n_{x\mu}\log(p\sigma_xL)-p\sigma_xL
+n'_x\log(\sigma_xL')-\sigma_xL'+\text{const}.
\]
若内部极值满足 \(0<\hat p<1\)，解得
\[
\hat\sigma_x=\frac{n'_x}{L'},
\qquad
\hat p=\frac{n_{x\mu}L'}{Ln'_x}.
\]
若该 \(\hat p>1\)，最大值会落在物理边界 \(p=1\)。

用误差传播近似，
\[
\hat p=\frac{L'}{L}\frac{n_{x\mu}}{n'_x},
\]
因此
\[
\sigma_{\hat p}\approx
\hat p\sqrt{\frac1{n_{x\mu}}+\frac1{n'_x}}.
\]
若以 \(n'_x\approx1000\) 作为与 \(n_x=10,L'=100L\) 相容的例子，则
\(\hat p=0.2\)，
\[
\sigma_{\hat p}\approx0.2\sqrt{\frac12+\frac1{1000}}=0.1416.
\]
这个无条件模型给出的近似区间与 9.5 的二项精确区间不同，因为它没有条件在固定的
\(n_x=10\) 上，而是用另一个实验估计高 \(x\) 总率。
\end{enumerate}
\end{solution}
\end{question}



\begin{question}
相互作用中产生的某粒子以相对于 \(z\) 的某一角度发射出来，如图 9.1
所示。探测器放在距离相互作用顶点为 \(d\) 处测量粒子垂直于 \(z\)
方向的位置 \(x\)。角度 \(\theta\) 定义为 \(z\) 轴与粒子径迹在 \((x,z)\)
平面上投影的夹角。假设测量值 \(x\) 可以看做以真值为中心值，标准偏差为
\(\sigma_x\) 的高斯变量。
\pyfig[0.68\linewidth]{figures/fig_09_07_detector_schematic.png}{图 9.1：探测器测量 \(x\) 并由几何关系确定散射角 \(\theta\) 的示意图。}


\begin{enumerate}
\item
  求 \(\cos\theta\) 置信水平为 \(1-\gamma\) 的中心置信区间。
\item
  取 \(d=1\,\mathrm{m}\)，\(\sigma_x=1\,\mathrm{mm}\)，并假设测量值为
  \(x=2.0\,\mathrm{mm}\)。求 \(\cos\theta\) 置信水平分别为 68.3\% 和
  95\% 的中心置信区间。
\end{enumerate}
\begin{solution}
几何关系为
\[
\tan\theta=\frac{x}{d},
\qquad
c\equiv\cos\theta=\frac{d}{\sqrt{d^2+x^2}}.
\]
测量值 \(\hat x\) 的中心置信区间为
\[
[x_-,x_+]=[\hat x-z_{1-\gamma/2}\sigma_x,
\hat x+z_{1-\gamma/2}\sigma_x].
\]
若整个区间在 \(x\ge0\) 或 \(x\le0\) 的同一侧，则 \(c(x)\) 随 \(|x|\)
单调下降，因此可直接把端点变换为
\[
c\in
\left[
\frac{d}{\sqrt{d^2+x_+^2}},
\frac{d}{\sqrt{d^2+x_-^2}}
\right]
\quad(x_->0).
\]
如果区间跨过 \(x=0\)，上端应取 \(c=1\)，下端由端点中较大的 \(|x|\) 给出。

题中
\[
d=1\,\mathrm m=1000\,\mathrm{mm},
\qquad
\sigma_x=1\,\mathrm{mm},
\qquad
\hat x=2.0\,\mathrm{mm}.
\]

\begin{itemize}
\item 68.3\% 中心区间近似取 \(z=1\)：\(x_-=1\,\mathrm{mm}\)，
\(x_+=3\,\mathrm{mm}\)。因此
\[
\cos\theta\in
[0.999995500,
0.999999500].
\]

\item 95\% 中心区间取 \(z=1.96\)：\(x_-=0.04\,\mathrm{mm}\)，
\(x_+=3.96\,\mathrm{mm}\)。因此
\[
\cos\theta\in
[0.999992159,
0.999999999].
\]
\end{itemize}
由于 \(x\ll d\)，角度很小，\(\cos\theta\) 的区间自然非常靠近 1。
\pyfig[0.72\linewidth]{figures/fig_09_07_cos_interval.png}{由 \(x\) 的中心置信区间通过单调变换得到 \(\cos\theta\) 的区间。}
\end{solution}
\end{question}

